\section{Experiments}
\label{sec:experiments}

\subsection{Protocol and statistics}
\label{sec:protocol}

\paragraph{Metrics.}
The headline metric is task-macro accuracy: the unweighted mean of the six task accuracies,
with failed requests counted as wrong. A Choice is correct if its argmax equals the gold
candidate; a Noul if $P(\text{true})>0.5$ agrees with the gold. NLL is the mean negative
log-likelihood of the gold candidate from a stable log-softmax. ECE uses ten equal-width bins
of the top candidate probability $p_{\max}$ (not the \texttt{confidence} field of
Section~\ref{sec:confidence}). ``Counterfactual all correct'' is the fraction of
counterfactual groups with a changing target in which every member is correct.

\paragraph{Intervals.}
All intervals are 95\% percentile intervals from 1,000 bootstrap resamples of components
(contexts, dialogues or tool situations). Differences between two models are paired: both
are evaluated on the same resampled components. Because the tool-call task has only six
development and three final components, its intervals understate the real uncertainty.

\paragraph{Development and final.}
The development split was observed repeatedly and was used to select the backbone, the arm
and every setting. The final split was opened once, after the release manifest was frozen
(Section~\ref{sec:final}). Development numbers are therefore optimistic in a way the final
number is not.

\subsection{Backbone selection}
\label{sec:leaderboard}

We compared eight pretrained backbones zero-shot through the same compiler and readout on the
development split (3,188 decisions, 1,161 components), with whole-payload encoding, no
calibration and no training. Seven ran with transformers 5.17.0 and Torch 2.12.1 in BF16 on
L40S GPUs (two GPUs for the 27B and 35B-A3B models). DeepSeek-V4.1-Flash ran with its official
reference inference code (tensor parallel 8, FP8 weights with FP4 experts) on eight B300
GPUs. Table~\ref{tab:leaderboard} and Figure~\ref{fig:leaderboard} give the results; per-task
accuracies are in Appendix~\ref{app:tables}.

\begin{table}[t]
\centering
\small
\caption{Zero-shot backbone comparison on the development split (3,188 decisions). Task
macro counts failed requests as wrong. Sizes are as recorded in our candidate records; ``--''
means not recorded there. DeepSeek-V4.1-Flash returned all-NaN logits for 36 rows
(1.1\%), which count as failures in the macro; its NLL and ECE use the 3,152 finite rows.}
\label{tab:leaderboard}
\setlength{\tabcolsep}{4pt}
\begin{adjustbox}{max width=\linewidth}
\begin{tabular}{@{}llcrrrr@{}}
\toprule
Backbone & Size & GPUs & Macro (\%) & NLL & ECE & CF all (\%) \\
\midrule
Qwen3.8-27B & 27.8B dense, hybrid & 2$\times$L40S & \textbf{85.8} & \textbf{0.511} & \textbf{0.033} & \textbf{78.4} \\
Gemma-4-12B-it & 12.0B dense & 1$\times$L40S & 84.9 & 1.350 & 0.130 & 72.2 \\
Qwen3.6-35B-A3B & 36B MoE, 3B active & 2$\times$L40S & 82.8 & 0.668 & 0.072 & 67.1 \\
Qwen3.5-9B & 9.7B dense, hybrid & 1$\times$L40S & 78.4 & 0.785 & 0.079 & 58.1 \\
DeepSeek-V4.1-Flash & 763B MoE, 8B active in prefill & 8$\times$B300 & 75.2 & 1.448 & 0.158 & 52.9 \\
A.X-4.0-Light & -- & 1$\times$L40S & 62.4 & 1.395 & 0.192 & 21.9 \\
Phi-4-mini-instruct & -- & 1$\times$L40S & 49.0 & 2.480 & 0.337 & 0.0 \\
Tri-7B & -- & 1$\times$L40S & 43.6 & 3.000 & 0.456 & 0.8 \\
\bottomrule
\end{tabular}
\end{adjustbox}
\end{table}

\begin{figure}[t]
\centering
\figfile[width=0.85\linewidth]{figures/leaderboard.pdf}
\caption{Zero-shot development task-macro accuracy of the eight backbones
(whole-payload encoding, failures counted wrong; ECE at the right), with \expone{} (piecewise
encoding, after training) and the piecewise zero-shot Qwen3.8-27B as reference markers; the
two encodings of the same zero-shot model differ slightly (85.8\% and 85.5\%). Development
split, 3,188 decisions; not the sealed final.}
\label{fig:leaderboard}
\end{figure}

\paragraph{Paired differences.}
On task macro, Qwen3.8-27B minus Gemma-4-12B-it is +0.91 points [$-0.69$, +2.25], so the two
are not distinguishable; Qwen3.8-27B exceeds Qwen3.6-35B-A3B by +3.01 [+1.37, +4.47] and
Qwen3.5-9B by +7.42 [+5.98, +9.03]. Qwen3.8-27B is also the best calibrated without training
(ECE 0.033), while Gemma-4-12B-it is overconfident (NLL 1.350).

\paragraph{DeepSeek-V4.1-Flash on a fair subset.}
Two factors penalize DeepSeek-V4.1-Flash in Table~\ref{tab:leaderboard}: NaN logits and
identifiers. On the 2,963 decisions with at most 62 candidates, where every model uses GLM's
identifiers, restricted to rows where DeepSeek's logits are finite, Qwen3.8-27B still leads
by +7.53 points [+5.94, +9.03]. DeepSeek needs eight GPUs and is not better, so its small
prefill-active parameter count does not help here; we keep it only as a possible reasoning
teacher. This conclusion applies to reading the first position without thinking, not to
DeepSeek's reasoning mode.

\begin{table}[t]
\centering
\small
\caption{Title selection accuracy (\%) by number of candidates $K$, development split,
75 questions per $K$ (DeepSeek-V4.1-Flash: finite rows only, 70 at $K=77$). Rows above the
rule use whole-payload encoding; the last two use piecewise encoding. For Qwen and Gemma,
$K\ge77$ uses extension letters as identifiers; the other models use GLM's multi-digit
numerals.}
\label{tab:kscaling}
\begin{tabular}{@{}lrrrr@{}}
\toprule
Model & $K=8$ & 77 & 200 & 255 \\
\midrule
Qwen3.8-27B & 92.0 & 76.0 & 70.7 & 66.7 \\
Gemma-4-12B-it & 93.3 & 74.7 & 72.0 & 68.0 \\
Qwen3.6-35B-A3B & 93.3 & 77.3 & 69.3 & 45.3 \\
Qwen3.5-9B & 89.3 & 81.3 & 41.3 & 34.7 \\
DeepSeek-V4.1-Flash & 94.7 & 2.9 & 0.0 & 0.0 \\
A.X-4.0-Light & 86.7 & 70.7 & 29.3 & 25.3 \\
Phi-4-mini-instruct & 68.0 & 13.3 & 9.3 & 8.0 \\
Tri-7B & 74.7 & 2.7 & 0.0 & 0.0 \\
\midrule
Qwen3.8-27B zero-shot, piecewise & 92.0 & 80.0 & 72.0 & 65.3 \\
\expone{} (dev), piecewise & 92.0 & 86.7 & 86.7 & 81.3 \\
\bottomrule
\end{tabular}
\end{table}

\paragraph{Many candidates and identifier choice.}
Table~\ref{tab:kscaling} and Figure~\ref{fig:kscaling} show accuracy as the number of
candidates grows. DeepSeek-V4.1-Flash has the best $K=8$ accuracy of all eight models
(94.7\%) and collapses at $K\ge77$ (2.9\%, 0.0\%, 0.0\%); Tri-7B and Phi-4-mini-instruct
collapse similarly. These three models used GLM's multi-digit numerals (\texttt{10},
\texttt{11}, \dots) as identifiers. The numerals are single tokens, but we suspect the models
tend to answer a number digit by digit, which a first-token readout cannot follow. Qwen and
Gemma, which used letters for $K\ge77$, degrade gradually. A.X-4.0-Light also used the
numerals throughout and falls more gradually (70.7\% at $K=77$, 25.3\% at $K=255$), so
candidate count itself is a second cause. The two effects are confounded in this design; an
identifier-only control has not been run. For future evaluations we use single-character
identifiers for every model.

\begin{figure}[t]
\centering
\figfile[width=0.75\linewidth]{figures/k_scaling.pdf}
\caption{Title selection accuracy against the number of candidates $K$ (development split,
75 questions per $K$). Shading marks $K\le62$, where all models share GLM's identifiers.
For $K\ge77$, filled markers denote models that use extension letters (Qwen, Gemma) and open
markers models that use GLM's multi-digit numerals (DeepSeek, A.X, Phi, Tri). Dashed lines use
piecewise encoding. DeepSeek-V4.1-Flash is scored over finite rows only.}
\label{fig:kscaling}
\end{figure}

\paragraph{Decision.}
We selected Qwen3.8-27B: it had the highest macro, the best zero-shot calibration and fits
on one GPU in BF16. Gemma-4-12B-it, statistically tied and less than half the size, remains
the main smaller candidate.

\subsection{LoRA fine-tuning and objective arms}
\label{sec:arms}

\paragraph{Setup.}
All arms start from the same backbone and piecewise inputs. The baseline is the same model
with the LoRA initialized at zero and evaluated without updates. The main arm,
\texttt{sft\_vocab}, trains LoRA with cross-entropy for one epoch: 1,592 steps of 32
questions (eight GPUs, four questions each, one question per forward), AdamW with learning
rate $10^{-4}$, 3\% warmup and cosine decay, gradient clipping at 1.0, on one node of eight
B300 GPUs (about 3.2\,s per step, peak memory 54.7\,GB per GPU). \texttt{sft\_pointer} adds the
zero-initialized pointer head of Section~\ref{sec:pointer}. The continuation arms
\texttt{rl}, \texttt{brier} and \texttt{ce} each continue from \texttt{sft\_vocab} for 300
steps at learning rate $2\times10^{-5}$ with the same seed; \texttt{rl} uses four samples per
question with a leave-one-out baseline. Temperatures are fitted on the calibration split and
applied to development. The monitor split (1,589 rows) follows the training distribution.
Appendix~\ref{app:hyper} lists all hyperparameters.

\begin{table}[t]
\centering
\small
\caption{Training arms on the development split (3,188 decisions, 1,161 components).
$\Delta$ is the paired difference in task macro against the zero-shot baseline, in points,
with a 95\% component-bootstrap interval. Tool call is the held-out task. ECE@$T$ is ECE after
the calibration-fitted temperature. Monitor is accuracy on the training-distribution monitor
split. K255 is title selection with 255 candidates (75 questions).}
\label{tab:arms}
\setlength{\tabcolsep}{3.5pt}
\begin{adjustbox}{max width=\linewidth}
\begin{tabular}{@{}lrlrrrrrrr@{}}
\toprule
Arm & Macro & $\Delta$ [95\%] & Tool & NLL & ECE & $T$ & ECE@$T$ & Monitor & K255 \\
\midrule
Zero-shot (same input) & 85.5 & -- & 81.8 & 0.526 & 0.033 & 1.26 & 0.045 & 83.1 & 65.3 \\
\texttt{sft\_vocab} (\expone) & \textbf{93.4} & +7.94 [+6.77, +9.02] & 95.0 & 0.236 & 0.015 & 1.15 & 0.011 & 95.5 & 81.3 \\
\texttt{sft\_pointer} & 92.9 & +7.42 [+6.05, +8.64] & 92.3 & 0.239 & 0.017 & 1.16 & 0.010 & 96.1 & 81.3 \\
\texttt{rl} (continued) & 93.3 & +7.80 [+6.60, +8.94] & 95.6 & 0.240 & 0.016 & 1.19 & 0.008 & 95.3 & 81.3 \\
\texttt{brier} (continued) & 93.1 & +7.67 [+6.50, +8.81] & 95.6 & 0.236 & 0.016 & 1.14 & 0.008 & 96.2 & 81.3 \\
\texttt{ce} (continued) & 93.2 & +7.68 [+6.48, +8.79] & 95.0 & 0.247 & 0.022 & 1.24 & 0.008 & 95.7 & 82.7 \\
\bottomrule
\end{tabular}
\end{adjustbox}
\end{table}

\paragraph{Supervised fine-tuning carries the gain.}
One epoch of LoRA SFT raises development task macro from 85.5\% to 93.4\% (+7.94 points
[+6.77, +9.02]) and counterfactual all-correct from 77.7\% to 94.7\%, and it lowers NLL from
0.526 to 0.236 and the monitor's Score normalized MAE from 0.134 to 0.089 (Table~\ref{tab:arms};
per-task results in Appendix~\ref{app:tables}). Restricted to the five trained tasks, the
gain is +6.87 points [+5.57, +8.11]. Calibration behaves differently before and after
training: the fitted temperature lowers the trained model's ECE from 0.015 to 0.011, but
raises the zero-shot model's from 0.033 to 0.045. Classification remains the weakest task
(77.6\%), with only 326 training rows and noisy labels.

\paragraph{The pointer head was not used.}
In \texttt{sft\_pointer} the pointer's scale started at zero and ended at $-0.0006$. LoRA solved
the task through the vocabulary readout, so the pointer received no consistent gradient.
This arm therefore did not test the pointer; Section~\ref{sec:limits} reports a retest.

\begin{table}[t]
\centering
\small
\caption{Paired differences between continuation arms on the development split (same 300
steps from \texttt{sft\_vocab}); 95\% component-bootstrap intervals.}
\label{tab:continuation}
\begin{tabular}{@{}lll@{}}
\toprule
Comparison & Accuracy (points) & NLL \\
\midrule
RL $-$ Brier & +0.12 [$-0.15$, +0.41] & +0.004 [0.000, +0.008] \\
RL $-$ CE & +0.11 [$-0.25$, +0.54] & $-0.007$ [$-0.013$, $-0.001$] \\
Brier $-$ CE & $-0.01$ [$-0.31$, +0.34] & $-0.011$ [$-0.017$, $-0.006$] \\
CE $-$ SFT & $-0.25$ [$-0.87$, +0.32] & +0.010 [+0.003, +0.018] \\
\bottomrule
\end{tabular}
\end{table}

\paragraph{Proper-score objectives protect probabilities, not accuracy.}
No continuation arm changes accuracy measurably (Table~\ref{tab:continuation}). Continuing
with cross-entropy makes the model overconfident (NLL +0.010 [+0.003, +0.018] over SFT). RL
and Brier avoid this: both have lower NLL than CE, and RL is slightly worse than Brier
(+0.004 [0.000, +0.008]), which is consistent with RL being a higher-variance estimator of the
same objective (Section~\ref{sec:objectives}). Where gold labels exist, direct Brier training
gives the same effect more cheaply and more stably than sampling.

\paragraph{Choice of candidate.}
We took \texttt{sft\_vocab}, the arm with the highest development macro and the simplest
design, attached its calibration temperature ($T=1.1489$) and named it \expone.

\subsection{Sealed final}
\label{sec:final}

\paragraph{Freeze.}
Before the final split was opened, a release manifest was frozen and committed. It binds the
base revision, the tokenizer and template hashes, the adapter's SHA256, the calibration
temperature 1.1489, the precision (BF16 weights, float32 readout), the decision rule (argmax,
ties by request order; Noul true only if $P(\text{true})>0.5$) and the hash of the final
file. The readout refuses to open the final split without such a manifest. The evaluation
path was identical to development (unmerged LoRA, batch size 1). The baseline was the same
piecewise input zero-shot, with the temperature 1.26 fitted on calibration beforehand. No
setting was changed after the final results were seen, and the final split will not be
reused for a later model. Appendix~\ref{app:compute} lists the manifest's hashes.

\begin{table}[t]
\centering
\small
\caption{Sealed final split, evaluated once (1,607 decisions, all scored, none failed).
$\Delta$ intervals are 95\% paired component-bootstrap intervals. The tool-call task was
never trained on and has only three final situations.}
\label{tab:final}
\begin{tabular}{@{}lrr@{}}
\toprule
Task (decisions) & \expone{} & Zero-shot (same input) \\
\midrule
Search passage selection (432) & 88.2 & 77.1 \\
Citation verification (275) & 100.0 & 91.6 \\
Tool-call review (96, held out) & 99.0 & 79.2 \\
External document screening (280) & 98.6 & 87.9 \\
Request routing (274) & 99.3 & 97.1 \\
Classification (250) & 82.8 & 74.0 \\
\midrule
Task macro & \textbf{94.6} & 84.5 \\
$\Delta$ task macro (points) & \multicolumn{2}{c}{+10.2 [+7.7, +12.3]} \\
$\Delta$ tool call (points) & \multicolumn{2}{c}{+19.8 [+12.5, +28.1]} \\
Counterfactual all correct & 93.6 & 74.0 \\
NLL, raw / at $T$ & 0.226 / 0.221 & 0.529 / 0.508 \\
ECE, raw / at $T$ & 0.016 / 0.011 & 0.029 / 0.030 \\
Temperature $T$ (fitted on calibration) & 1.1489 & 1.2596 \\
\bottomrule
\end{tabular}
\end{table}

\paragraph{Results.}
Table~\ref{tab:final} and Figure~\ref{fig:final} give the result. \expone{} reaches 94.6\%
task macro against 84.5\% for the same-input zero-shot model, a paired difference of +10.2
points [+7.7, +12.3]. The final result is not lower than development (93.4\%). The
calibrated ECE is 0.011. All 1,607 requests were scored. Within search passage selection,
title selection weakens as candidates increase: 94.7\% at $K=8$, 78.9\% at 77, 71.1\% at 200
and 73.7\% at 255 (the zero-shot model ranges from 47\% to 71\% across these four settings).
Classification is the weakest task at 82.8\%.

\paragraph{What this does not establish.}
The final labels have not been reviewed by a person. The project's release gate G1 requires
non-inferiority against the GLM-5.3-Flash direct-logit readout on the same evaluation, which
has not been run; this result alone therefore does not decide any release gate. The final
split is Korean only.

\begin{figure}[t]
\centering
\figfile[width=0.85\linewidth]{figures/final_tasks.pdf}
\caption{Sealed final accuracy per task for \expone{} and the same-input zero-shot model,
with the number of decisions per task. Tool-call review was never trained on.}
\label{fig:final}
\end{figure}

\subsection{The held-out tool-call task}
\label{sec:heldout}

Tool-call review was absent from training. On development its accuracy rose from 81.8\% to
95.0\% (+13.3 points [+11.6, +14.8], six situations, 181 decisions) and on the sealed final
from 79.2\% to 99.0\% (+19.8 [+12.5, +28.1], three situations, 96 decisions). This is our
clearest generalization signal, because the other five tasks come from generators that also
produced training data. It rests on few independent units, however: with six and three
situations, the component bootstrap cannot reflect situation-level variability well.

\subsection{Ten unseen tool situations}
\label{sec:toolood}

To test the tool-call result on more situations, we wrote a catalog of ten situations that
no training or evaluation split uses: cron jobs, DNS, payment refunds, training checkpoints,
Linux accounts, CI, newsletters, a search index, backups and web deployment. Rows were built
with the evaluation's generator (311 rows). The situations and labels were written by an AI
assistant (Claude) following the conventions of the original catalog and have not been
reviewed by a person.

\begin{table}[t]
\centering
\small
\caption{Ten unseen tool-call situations (311 rows; AI-written labels, not human-reviewed).
No intervals were computed for this set.}
\label{tab:toolood}
\begin{tabular}{@{}lrr@{}}
\toprule
Question family (rows) & \expone{} & Zero-shot \\
\midrule
Plan mismatch (160) & 95.6 & 70.6 \\
Irreversible (80) & 98.8 & 90.0 \\
Scope (71) & \textbf{88.7} & \textbf{91.5} \\
\midrule
Family macro & 94.4 & 84.1 \\
Counterfactual all correct & 87.4 & 53.2 \\
\bottomrule
\end{tabular}
\end{table}

The family macro of 94.4\% (Table~\ref{tab:toolood}) is close to the development tool-call
accuracy (95.0\%), so the transfer holds on new situations. The scope family is the
exception: \expone{} is 2.8 points below the untrained model (88.7\% against 91.5\%, 71 rows,
so the difference is within noise at this size). The ability to compare an action with a
stated plan appears to transfer from the trained tasks, while judging whether an action
exceeds the request may have been disturbed.

\subsection{Author-written bilingual probes}
\label{sec:probes}

We wrote 108 requests (126 questions) in situations the evaluation generators never
produced: customer-support routing, English tool-call review, citations involving arithmetic
and negation, reviews and e-mails with hidden instructions to an AI, human-facing notices
and quotations that describe attacks, model-tier routing, Score rubrics (answer quality,
code-review severity), spam, unsubscribe requests and satire, policy violations, code
defects, three-question requests, and known weak spots (sums, elapsed dates, thresholds,
units). Each concept has a Korean and an English version with the same expected answer
(63 questions per language). The probes and their labels were written by an AI assistant
(Claude) and have not been reviewed by a person; the set is small and has many easy items, so
it has a ceiling. It is a sanity check, not a benchmark, and was never used for training,
selection or calibration.

\begin{table}[t]
\centering
\small
\caption{Author-written bilingual probes (126 questions, 63 per language; AI-written labels,
not human-reviewed).}
\label{tab:probes}
\begin{tabular}{@{}lrr@{}}
\toprule
 & \expone{} & Zero-shot \\
\midrule
All & 98.4 (124/126) & 95.2 (120/126) \\
Korean / English & 98.4 / 98.4 & 93.7 / 96.8 \\
Korean and English pair gives the same answer & 96.8 & 92.1 \\
Mean probability on the expected answer & 0.959 & 0.877 \\
\bottomrule
\end{tabular}
\end{table}

\expone{} missed two questions (Table~\ref{tab:probes}), both sums: it accepted an English
invoice total that was wrong ($P(\text{yes})=0.56$) and rejected a Korean total that was right
($P(\text{yes})=0.46$). Both probabilities are near 0.5, so the model was wrong without being
confident. The untrained model missed six: the same two sums; ``540 of 1,200 people'' judged
as supporting ``more than half'' in both languages (with probabilities of 0.68 to 0.70); 2,300\,g
judged to be at most 2\,kg; and a Docker image prune judged irreversible. Among the probes both
models answered correctly: a satirical review read as negative, a hidden
\texttt{display:none} instruction flagged as injection while a human-facing administrator
notice and a quotation describing an attack were not, a cancellation request that looked
like a billing question routed to cancellation, \texttt{terraform apply} judged to exceed a
request for a preview, and a full card number judged a policy violation. Arithmetic
remains a weakness, and, as our specification requires, it belongs in code.

\subsection{Earlier objective comparison on the GLM reference}
\label{sec:glmrl}

Before the student direction, we compared three objectives on GLM-5.3-Flash with a LoRA
adapter, from the same parent checkpoint, for 35 updates each with the same data order and
learning rate, on a different development set of 1,024 components (896 categorical
decisions). Table~\ref{tab:glmrl} summarizes the result.

\begin{table}[t]
\centering
\small
\caption{Earlier three-objective comparison on GLM-5.3-Flash (single seed; development data of
a different evaluation: 896 categorical decisions, 562 of them native Korean).}
\label{tab:glmrl}
\setlength{\tabcolsep}{4pt}
\begin{adjustbox}{max width=\linewidth}
\begin{tabular}{@{}lrrrr@{}}
\toprule
 & Parent & Proper-score RL & Direct Brier & Correctness-only RL \\
\midrule
Correct / 896 & 796 & 802 & 802 & 799 \\
Native Korean correct / 562 & 509 & 509 & 510 & 508 \\
NLL & 0.34126 & 0.33594 & 0.33365 & 0.34545 \\
ECE & 0.04515 & 0.04344 & 0.04163 & 0.05187 \\
\bottomrule
\end{tabular}
\end{adjustbox}
\end{table}

Proper-score RL changed accuracy by +0.67 points [0.00, +1.45] over the parent, and RL minus
Brier was 0.00 points [$-0.67$, +0.67]. Temperature scaling alone applied to the parent
(development cross-fitting) gave lower NLL (0.32550) and ECE (0.02061) than raw RL. We
concluded that with gold labels, a direct proper-score loss plus temperature is the default,
and that RL should return only where an outcome without gold labels can serve as the reward.
The student experiments of Section~\ref{sec:arms} reproduce the same pattern.

Two further observations from this run matter for how RL is used. First, most sampled
updates carried no RL signal: in 77.5\% of the proper-score RL components and 81.4\% of the
correctness-only components, all four samples received the same reward and the leave-one-out
advantage was zero,
so only the retention term moved the parameters. Second, rewarding correctness alone is not a
neutral simplification for a model whose product is a probability. It was the only arm whose
NLL and ECE rose above the parent's (Table~\ref{tab:glmrl}), which is what an improper reward
that pays for confident correct answers should do.

\begin{figure}[t]
\centering
\figfile{figures/rl.pdf}
\caption{The two reinforcement-learning experiments. Top: continuing \expone{} for 300 steps
(Table~\ref{tab:arms}). Bottom: GLM-5.3-Flash, 35 updates per objective from one parent
(Table~\ref{tab:glmrl}). With gold labels, proper-score RL tracks the Brier loss it estimates;
correctness-only reward is the arm that hurts calibration.}
\label{fig:rl}
\end{figure}

\paragraph{Where RL returns.}
Proposition~1 assumes that the loss can be evaluated for every candidate, which is true
exactly when the gold label is known. When the only signal is a downstream outcome, such as a
workflow that succeeded or a reviewer who overturned a decision, there is no closed-form loss to
compute, and a sampled proper-score reward is the natural estimator. That is the setting of our
next RL experiment, to be compared against the supervised and temperature-calibrated
baseline of this section rather than against the untrained model.
